%% ============================================================
%%  Chapter 4: Experimental Evaluation
%% ============================================================
\chapter{Experimental Evaluation}
\label{chap:experiments}

% GUIDANCE: 15--20 pages. This is your results chapter.
% Write the setup sections NOW (they don't need results).
% Leave \todo placeholders for figures and tables.
% Fill in numbers once experiments finish.


% ── 4.1 ─────────────────────────────────────────────────────────
\section{Experimental Setup}
\label{sec:exp_setup}


\subsection{Environment Configuration}
\label{subsec:exp_environment}

All experiments are run in the same warehouse gridworld simulation stack,
with map-specific geometry, agent count, induct stations, and eject stations
loaded from YAML configuration files. One simulation timestep corresponds to
one discrete planning/execution update in the environment. Unless otherwise
noted, all allocation methods share the same downstream path-planning and
collision-avoidance pipeline to isolate the effect of task-allocation logic.


\subsection{Parameters Swept}
\label{subsec:exp_parameters}

All experiments sweep the following parameters:

\begin{itemize}
  \item \textbf{Communication range} is map-scaled from diagonal fractions
    $\{0.1, 0.2, 0.35, 0.6, 1.2\}$ and rounded to grid units.
  \item \textbf{Batch tasks per induct station} in full mode:
    $\{8, 10, 12, 15\}$ (reliability-first sweep).
  \item \textbf{Batch tasks per induct station} in stress mode:
    $\{10, 20, 30, 40\}$ (overload analysis, reported separately).
  \item \textbf{Steady-state arrival fractions} in full mode:
    $\{0.25, 0.4, 0.55, 0.7, 0.85\}$ times the map-specific
    estimated steady-state capacity.
  \item \textbf{Steady-state arrival fractions} in quick mode:
    $\{0.5, 0.8\}$ times the map-specific estimated steady-state capacity.
  \item \textbf{Steady-state arrival fractions} in stress mode:
    $\{0.6, 0.8, 1.0, 1.2, 1.4\}$ times the map-specific estimated
    steady-state capacity.
  \item \textbf{Rerun interval} (Dynamic GCBBA only)
        $\in \{10, 25, 50, 100, 200\}$ timesteps, with $ri = 50$
        as the canonical configuration.
\end{itemize}

Each configuration is repeated across two random seeds ($42$, $123$) in full mode.
The maximum timestep budget is 1500 for steady-state and 3000 for batch runs.

For steady-state experiments, quick mode is now a low-load smoke check,
full mode is intentionally conservative so that the main thesis figures focus on
stable operating regions, and stress mode is reserved for overload behavior.

A complete sweep-count table is omitted here for brevity and will be added in
the final manuscript revision after steady-state result integration.


\subsection{Batch Reliability Screening Across Maps}
\label{subsec:exp_batch_reliability_screen}

To ensure that the main thesis comparisons are based on stable operating regimes,
we screened each map using the latest full-mode batch runs and marked a load point
as \emph{reliable} when the wall-clock timeout rate was at most $20\%$.

Table~\ref{tab:batch_reliability_screen} summarizes the maximum reliable batch load
per map and the recommended reporting role in this thesis.

\begin{table}[htbp]
  \centering
  \small
  \caption{Batch reliability screen from full-mode runs. A map is marked
       as main-analysis-ready when it provides usable low-to-mid load
       comparisons without timeout-dominated behavior.}
  \label{tab:batch_reliability_screen}
  \begin{tabular}{lccc}
    \hline
    Map & Max reliable task count & Completed-only makespan usable & Recommended use \\
    \hline
    gridworld\_crossdock       & 90  & Yes & Main analysis \\
    gridworld\_warehouse\_small & 120 & Partial (best at lower loads) & Main analysis \\
    gridworld\_warehouse\_large & 144 & Partial (exclude top load in main) & Main analysis \\
    gridworld\_kiva            & 75  & No (timestep-ceiling dominated) & Secondary / partial-completion \\
    gridworld\_kiva\_large      & N/A & No & Stress / appendix \\
    gridworld\_shelf\_aisle      & N/A & No & Stress / appendix \\
    \hline
  \end{tabular}
\end{table}

The resulting presentation policy is: (i) use crossdock, warehouse\_small,
and warehouse\_large as the primary cross-map batch comparison set,
(ii) report kiva primarily through completion-fraction and timeout behavior,
and (iii) treat kiva\_large and shelf\_aisle as stress-test regimes.

The compact batch figures used for this screening were generated from:
\begin{itemize}
  \item \texttt{results/experiments/gridworld\_warehouse\_small/20260425\_012846/batch\_results/compact\_views/}
  \item \texttt{results/experiments/gridworld\_crossdock/20260425\_022224/batch\_results/compact\_views/}
  \item \texttt{results/experiments/gridworld\_kiva/20260425\_030403/batch\_results/compact\_views/}
  \item \texttt{results/experiments/gridworld\_warehouse\_large/20260425\_044219/batch\_results/compact\_views/}
  \item \texttt{results/experiments/gridworld\_kiva\_large/20260425\_062503/batch\_results/compact\_views/}
  \item \texttt{results/experiments/gridworld\_shelf\_aisle/20260425\_074507/batch\_results/compact\_views/}
\end{itemize}

For each run, the primary reliability visuals are:
\texttt{compact\_heatmap\_completion\_rate.png},
\texttt{compact\_heatmap\_timeout\_rate.png},
\texttt{compact\_capacity\_curve.png}, and
\texttt{compact\_completion\_degradation\_by\_comm\_range.png}.
When available, completed-only makespan is shown via
\texttt{compact\_makespan\_completed\_by\_comm\_range.png}.

\subsubsection*{Main-Map Reliability Figures}

\begin{figure}[htbp]
  \centering
  \begin{subfigure}[t]{0.32\textwidth}
    \centering
    \includegraphics[width=\textwidth]{../results/experiments/gridworld_crossdock/20260425_022224/batch_results/compact_views/compact_heatmap_completion_rate.png}
    \caption{Crossdock}
  \end{subfigure}
  \hfill
  \begin{subfigure}[t]{0.32\textwidth}
    \centering
    \includegraphics[width=\textwidth]{../results/experiments/gridworld_warehouse_small/20260425_012846/batch_results/compact_views/compact_heatmap_completion_rate.png}
    \caption{Warehouse-small}
  \end{subfigure}
  \hfill
  \begin{subfigure}[t]{0.32\textwidth}
    \centering
    \includegraphics[width=\textwidth]{../results/experiments/gridworld_warehouse_large/20260425_044219/batch_results/compact_views/compact_heatmap_completion_rate.png}
    \caption{Warehouse-large}
  \end{subfigure}
  \caption{Completion-rate heatmaps for the three main-analysis maps used in the batch reliability screen.}
  \label{fig:batch_completion_heatmaps_main_maps}
\end{figure}

\begin{figure}[htbp]
  \centering
  \begin{subfigure}[t]{0.32\textwidth}
    \centering
    \includegraphics[width=\textwidth]{../results/experiments/gridworld_crossdock/20260425_022224/batch_results/compact_views/compact_heatmap_timeout_rate.png}
    \caption{Crossdock}
  \end{subfigure}
  \hfill
  \begin{subfigure}[t]{0.32\textwidth}
    \centering
    \includegraphics[width=\textwidth]{../results/experiments/gridworld_warehouse_small/20260425_012846/batch_results/compact_views/compact_heatmap_timeout_rate.png}
    \caption{Warehouse-small}
  \end{subfigure}
  \hfill
  \begin{subfigure}[t]{0.32\textwidth}
    \centering
    \includegraphics[width=\textwidth]{../results/experiments/gridworld_warehouse_large/20260425_044219/batch_results/compact_views/compact_heatmap_timeout_rate.png}
    \caption{Warehouse-large}
  \end{subfigure}
  \caption{Timeout-rate heatmaps for main-analysis maps. Low-to-mid loads are stable, while top loads show map-dependent degradation.}
  \label{fig:batch_timeout_heatmaps_main_maps}
\end{figure}

\begin{figure}[htbp]
  \centering
  \begin{subfigure}[t]{0.32\textwidth}
    \centering
    \includegraphics[width=\textwidth]{../results/experiments/gridworld_crossdock/20260425_022224/batch_results/compact_views/compact_capacity_curve.png}
    \caption{Crossdock}
  \end{subfigure}
  \hfill
  \begin{subfigure}[t]{0.32\textwidth}
    \centering
    \includegraphics[width=\textwidth]{../results/experiments/gridworld_warehouse_small/20260425_012846/batch_results/compact_views/compact_capacity_curve.png}
    \caption{Warehouse-small}
  \end{subfigure}
  \hfill
  \begin{subfigure}[t]{0.32\textwidth}
    \centering
    \includegraphics[width=\textwidth]{../results/experiments/gridworld_warehouse_large/20260425_044219/batch_results/compact_views/compact_capacity_curve.png}
    \caption{Warehouse-large}
  \end{subfigure}
  \caption{Critical-load frontier curves (completion threshold $\geq 0.95$) for main-analysis maps.}
  \label{fig:batch_capacity_main_maps}
\end{figure}

\begin{figure}[htbp]
  \centering
  \begin{subfigure}[t]{0.32\textwidth}
    \centering
    \includegraphics[width=\textwidth]{../results/experiments/gridworld_crossdock/20260425_022224/batch_results/compact_views/compact_completion_degradation_by_comm_range.png}
    \caption{Crossdock}
  \end{subfigure}
  \hfill
  \begin{subfigure}[t]{0.32\textwidth}
    \centering
    \includegraphics[width=\textwidth]{../results/experiments/gridworld_warehouse_small/20260425_012846/batch_results/compact_views/compact_completion_degradation_by_comm_range.png}
    \caption{Warehouse-small}
  \end{subfigure}
  \hfill
  \begin{subfigure}[t]{0.32\textwidth}
    \centering
    \includegraphics[width=\textwidth]{../results/experiments/gridworld_warehouse_large/20260425_044219/batch_results/compact_views/compact_completion_degradation_by_comm_range.png}
    \caption{Warehouse-large}
  \end{subfigure}
  \caption{Completion-vs-load degradation curves by communication range for main-analysis maps. These plots show graceful vs sharp degradation regimes and support the load-screening decisions in Table~\ref{tab:batch_reliability_screen}.}
  \label{fig:batch_degradation_main_maps}
\end{figure}

\begin{figure}[htbp]
  \centering
  \begin{subfigure}[t]{0.32\textwidth}
    \centering
    \includegraphics[width=\textwidth]{../results/experiments/gridworld_crossdock/20260425_022224/batch_results/compact_views/compact_tradeoff.png}
    \caption{Crossdock}
  \end{subfigure}
  \hfill
  \begin{subfigure}[t]{0.32\textwidth}
    \centering
    \includegraphics[width=\textwidth]{../results/experiments/gridworld_warehouse_small/20260425_012846/batch_results/compact_views/compact_tradeoff.png}
    \caption{Warehouse-small}
  \end{subfigure}
  \hfill
  \begin{subfigure}[t]{0.32\textwidth}
    \centering
    \includegraphics[width=\textwidth]{../results/experiments/gridworld_warehouse_large/20260425_044219/batch_results/compact_views/compact_tradeoff.png}
    \caption{Warehouse-large}
  \end{subfigure}
  \caption{Completion-versus-wall-time tradeoff for main-analysis maps. Pareto-favorable regions correspond to high completion at lower wall-time.}
  \label{fig:batch_tradeoff_main_maps}
\end{figure}

\subsubsection*{Stress-Regime Figures}

\begin{figure}[htbp]
  \centering
  \begin{subfigure}[t]{0.32\textwidth}
    \centering
    \includegraphics[width=\textwidth]{../results/experiments/gridworld_kiva/20260425_030403/batch_results/compact_views/compact_heatmap_timeout_rate.png}
    \caption{Kiva (secondary)}
  \end{subfigure}
  \hfill
  \begin{subfigure}[t]{0.32\textwidth}
    \centering
    \includegraphics[width=\textwidth]{../results/experiments/gridworld_kiva_large/20260425_062503/batch_results/compact_views/compact_heatmap_timeout_rate.png}
    \caption{Kiva-large (stress)}
  \end{subfigure}
  \hfill
  \begin{subfigure}[t]{0.32\textwidth}
    \centering
    \includegraphics[width=\textwidth]{../results/experiments/gridworld_shelf_aisle/20260425_074507/batch_results/compact_views/compact_heatmap_timeout_rate.png}
    \caption{Shelf-aisle (stress)}
  \end{subfigure}
  \caption{Timeout-rate heatmaps for secondary/stress maps. These are interpreted as robustness stress regimes rather than primary efficiency comparisons.}
  \label{fig:batch_timeout_heatmaps_stress_maps}
\end{figure}

\subsubsection*{Steady-State Figures (Rolling Integration)}

Steady-state results are being integrated incrementally as map runs finish.
At this revision, \texttt{gridworld\_warehouse\_small} and
\texttt{gridworld\_crossdock} are available as the current stable-regime SS maps,
and \texttt{gridworld\_kiva} is now available as a stress-regime SS map.
The SS thesis presentation uses the same compact-view structure across maps:
throughput, wait-time, capacity frontier, and completion progress split by stop
reason.

The newest steady-state runs are more representative than the earlier warm-started
versions: warehouse-small now completes more cleanly with much lower wait time,
while crossdock still shows timeout pressure but with substantially less wall-clock
domination than before. Kiva remains ceiling-dominated, so it is interpreted as a
stress-regime robustness case rather than a stable-regime efficiency comparison.
The saturation early-stop did not trigger in these maps, so the current
improvement comes from the revised arrival policy rather than from aggressive
early termination.

Current steady-state compact-view path:
\begin{itemize}
  \item \texttt{results/experiments/gridworld\_warehouse\_small/20260426\_135556/ss\_results/compact\_views/}
  \item \texttt{results/experiments/gridworld\_crossdock/20260426\_141042/ss\_results/compact\_views/}
  \item \texttt{results/experiments/gridworld\_kiva/20260426\_195056/ss\_results/compact\_views/}
\end{itemize}

\begin{figure}[htbp]
  \centering
  \begin{subfigure}[t]{0.48\textwidth}
    \centering
    \includegraphics[width=\textwidth]{../results/experiments/gridworld_warehouse_small/20260426_135556/ss_results/compact_views/compact_heatmap_throughput.png}
    \caption{Warehouse-small: throughput}
  \end{subfigure}
  \hfill
  \begin{subfigure}[t]{0.48\textwidth}
    \centering
    \includegraphics[width=\textwidth]{../results/experiments/gridworld_warehouse_small/20260426_135556/ss_results/compact_views/compact_heatmap_wait_time.png}
    \caption{Warehouse-small: wait time}
  \end{subfigure}

  \vspace{0.7em}

  \begin{subfigure}[t]{0.48\textwidth}
    \centering
    \includegraphics[width=\textwidth]{../results/experiments/gridworld_crossdock/20260426_141042/ss_results/compact_views/compact_heatmap_throughput.png}
    \caption{Crossdock: throughput}
  \end{subfigure}
  \hfill
  \begin{subfigure}[t]{0.48\textwidth}
    \centering
    \includegraphics[width=\textwidth]{../results/experiments/gridworld_crossdock/20260426_141042/ss_results/compact_views/compact_heatmap_wait_time.png}
    \caption{Crossdock: wait time}
  \end{subfigure}
  \caption{Steady-state throughput and wait-time heatmaps for the currently completed SS maps.}
  \label{fig:ss_throughput_wait_heatmaps_current_maps}
\end{figure}

\begin{figure}[htbp]
  \centering
  \begin{subfigure}[t]{0.48\textwidth}
    \centering
    \includegraphics[width=\textwidth]{../results/experiments/gridworld_warehouse_small/20260426_135556/ss_results/compact_views/compact_capacity_curve.png}
    \caption{Warehouse-small}
  \end{subfigure}
  \hfill
  \begin{subfigure}[t]{0.48\textwidth}
    \centering
    \includegraphics[width=\textwidth]{../results/experiments/gridworld_crossdock/20260426_141042/ss_results/compact_views/compact_capacity_curve.png}
    \caption{Crossdock}
  \end{subfigure}
  \caption{Steady-state capacity frontier for the currently completed SS maps.}
  \label{fig:ss_capacity_current_maps}
\end{figure}

\begin{figure}[htbp]
  \centering
  \begin{subfigure}[t]{0.48\textwidth}
    \centering
    \includegraphics[width=\textwidth]{../results/experiments/gridworld_warehouse_small/20260426_135556/ss_results/compact_views/compact_completion_progress_all_tasks_completed.png}
    \caption{Warehouse-small}
  \end{subfigure}
  \hfill
  \begin{subfigure}[t]{0.48\textwidth}
    \centering
    \includegraphics[width=\textwidth]{../results/experiments/gridworld_crossdock/20260426_141042/ss_results/compact_views/compact_completion_progress_all_tasks_completed.png}
    \caption{Crossdock}
  \end{subfigure}
  \caption{Steady-state completion progress for runs that finished all tasks.}
  \label{fig:ss_completion_progress_completed_current_maps}
\end{figure}

\begin{figure}[htbp]
  \centering
  \begin{subfigure}[t]{0.48\textwidth}
    \centering
    \includegraphics[width=\textwidth]{../results/experiments/gridworld_warehouse_small/20260426_135556/ss_results/compact_views/compact_completion_progress_timestep_ceiling.png}
    \caption{Warehouse-small}
  \end{subfigure}
  \hfill
  \begin{subfigure}[t]{0.48\textwidth}
    \centering
    \includegraphics[width=\textwidth]{../results/experiments/gridworld_crossdock/20260426_141042/ss_results/compact_views/compact_completion_progress_timestep_ceiling.png}
    \caption{Crossdock}
  \end{subfigure}
  \caption{Steady-state completion progress for runs that reached the timestep ceiling.}
  \label{fig:ss_completion_progress_timestep_ceiling_current_maps}
\end{figure}

\paragraph{Kiva Stress-Regime Steady-State Insertion.}
The latest Kiva steady-state run is included as a stress-regime result.
Kiva does not currently produce a completed-runs progress panel in this run
because the completed stop reason is absent; therefore, the chapter reports
throughput/wait heatmaps, capacity frontier, and timestep-ceiling completion progress.
This run is useful for robustness interpretation, but it should not be read as a
real-time control benchmark.

\begin{figure}[htbp]
  \centering
  \begin{subfigure}[t]{0.48\textwidth}
    \centering
    \includegraphics[width=\textwidth]{../results/experiments/gridworld_kiva/20260426_195056/ss_results/compact_views/compact_heatmap_throughput.png}
    \caption{Kiva: throughput}
  \end{subfigure}
  \hfill
  \begin{subfigure}[t]{0.48\textwidth}
    \centering
    \includegraphics[width=\textwidth]{../results/experiments/gridworld_kiva/20260426_195056/ss_results/compact_views/compact_heatmap_wait_time.png}
    \caption{Kiva: wait time}
  \end{subfigure}
  \caption{Steady-state throughput and wait-time heatmaps for Kiva (stress regime).}
  \label{fig:ss_kiva_heatmaps_stress}
\end{figure}

\begin{figure}[htbp]
  \centering
  \begin{subfigure}[t]{0.48\textwidth}
    \centering
    \includegraphics[width=\textwidth]{../results/experiments/gridworld_kiva/20260426_195056/ss_results/compact_views/compact_capacity_curve.png}
    \caption{Kiva: capacity frontier}
  \end{subfigure}
  \hfill
  \begin{subfigure}[t]{0.48\textwidth}
    \centering
    \includegraphics[width=\textwidth]{../results/experiments/gridworld_kiva/20260426_195056/ss_results/compact_views/compact_completion_progress_timestep_ceiling.png}
    \caption{Kiva: timestep-ceiling completion progress}
  \end{subfigure}
  \caption{Kiva stress-regime steady-state summary: capacity frontier and completion progress conditioned on timestep-ceiling stops.}
  \label{fig:ss_kiva_capacity_progress_stress}
\end{figure}

\paragraph{Batch-Style Multi-Map SS Template (to be filled as runs finish).}
The final steady-state section will keep the same map grouping policy as batch
(main-analysis maps, secondary maps, stress-regime maps) after SS reliability
screening is complete. For each map, the chapter will include the following files
from \texttt{.../ss\_results/compact\_views/}:
\begin{itemize}
  \item \texttt{compact\_heatmap\_throughput.png}
  \item \texttt{compact\_heatmap\_wait\_time.png}
  \item \texttt{compact\_capacity\_curve.png}
  \item \texttt{compact\_completion\_progress\_all\_tasks\_completed.png}
  \item \texttt{compact\_completion\_progress\_timestep\_ceiling.png}
\end{itemize}

Pending steady-state map folders (currently running):
\begin{itemize}
  \item \texttt{results/experiments/gridworld\_warehouse\_large/<ss\_timestamp>/ss\_results/compact\_views/}
  \item \texttt{results/experiments/gridworld\_kiva\_large/<ss\_timestamp>/ss\_results/compact\_views/}
  \item \texttt{results/experiments/gridworld\_shelf\_aisle/<ss\_timestamp>/ss\_results/compact\_views/}
\end{itemize}


\subsection{Metrics}
\label{subsec:exp_metrics}

The primary batch metrics used in this chapter are:
\begin{itemize}
  \item \textbf{Completion rate}: fraction of tasks completed within run limits.
  \item \textbf{Timeout rate}: fraction of runs that hit the wall-clock ceiling.
  \item \textbf{Completed-only makespan}: completion time, reported only for fully completed runs.
  \item \textbf{Wall-clock runtime}: end-to-end runtime per run.
  \item \textbf{Critical-load frontier}: maximum load with completion at or above 95\%.
\end{itemize}

The primary steady-state metrics used in this chapter are:
\begin{itemize}
  \item \textbf{Throughput}: completed steady-state tasks per timestep.
  \item \textbf{Average task wait time}: responsiveness under continuous arrivals.
  \item \textbf{Queue depth and queue-pressure indicators}: stability under sustained load.
  \item \textbf{Arrival-rate capacity frontier}: highest arrival rate sustaining near-peak throughput.
  \item \textbf{Throughput--wait tradeoff}: efficiency vs latency operating points.
\end{itemize}

Secondary diagnostics (idle ratio, balance, deadlocks, and allocation timing)
are available in the compact and detailed suites and are used where relevant in
the appendix-level discussion.


\subsection{Baselines}
\label{subsec:exp_baselines}

The compared methods are LCBA (internal key \texttt{gcbba}), CBBA, SGA,
and DMCHBA. All methods use the same simulation and path-planning stack so
that differences primarily reflect allocation behavior. DMCHBA is retained as
a high-performance reference baseline, while LCBA/CBBA/SGA provide
decentralized-style comparisons with varying robustness under communication and
load stress.


% ── 4.2 ─────────────────────────────────────────────────────────
\section{Communication Robustness}
\label{sec:exp_comm_robustness}

% THE KEY RESULT. Dynamic GCBBA completes tasks where static fails.

\subsection{Task Completion Rate}
\label{subsec:exp_completion_rate}

This subsection will be finalized after steady-state outputs are available.
For the current batch-only reliability analysis, completion behavior is
summarized in Figure~\ref{fig:batch_completion_heatmaps_main_maps} and
Figure~\ref{fig:batch_degradation_main_maps}.

\begin{figure}[htbp]
  \centering
  % \includegraphics[width=\textwidth]{completion_rate}
  \emph{[Steady-state completion-rate figure pending SS run outputs.]}
  \caption{Task completion rate by communication range, aggregated
           across all task loads and seeds. The shaded region
           indicates disconnected graph configurations
           ($\comm < 13$).}
  \label{fig:completion_rate}
\end{figure}

\begin{figure}[htbp]
  \centering
  % \includegraphics[width=\textwidth]{completion_rate_by_tpi}
  \emph{[Steady-state completion-by-load figure pending SS run outputs.]}
  \caption{Task completion rate broken down by task load.}
  \label{fig:completion_rate_by_tpi}
\end{figure}


\subsection{Makespan Under Varying Connectivity}
\label{subsec:exp_makespan_comm}

This subsection is reserved for the steady-state analysis pass.
For current batch data, makespan is interpreted only on completed runs and is
used as a secondary metric behind completion and timeout.

\begin{figure}[htbp]
  \centering
  % \includegraphics[width=\textwidth]{makespan_vs_comm_range}
  \emph{[Steady-state makespan-vs-communication figure pending SS outputs.]}
  \caption{Makespan vs.\ communication range for all methods
           across three task loads.}
  \label{fig:makespan_vs_comm}
\end{figure}


% ── 4.3 ─────────────────────────────────────────────────────────
\section{Optimality Analysis}
\label{sec:exp_optimality}

% Second key result: GCBBA stays near SGA on connected graphs.

\subsection{Sub-Optimality Ratio vs.\ SGA}
\label{subsec:exp_suboptimality}

This sub-optimality analysis is deferred to the steady-state section once the
SS runs are complete.

\begin{figure}[htbp]
  \centering
  % \includegraphics[width=\textwidth]{optimality_ratio}
  \emph{[Sub-optimality ratio figure pending SS outputs.]}
  \caption{Sub-optimality ratio (method makespan / SGA makespan)
           vs.\ communication range. Only runs where both the
           method and SGA completed are included.
           The dashed line at 2.0 indicates the theoretical
           worst-case guarantee for GCBBA on connected graphs.}
  \label{fig:optimality_ratio}
\end{figure}


\subsection{Makespan Gap Relative to SGA}
\label{subsec:exp_makespan_gap}

This subsection remains pending steady-state outputs.

\begin{figure}[htbp]
  \centering
  % \includegraphics[width=\textwidth]{improvement_ratio}
  \emph{[Improvement-ratio figure pending SS outputs.]}
  \caption{Makespan gap relative to SGA, expressed as a
           percentage. Positive values indicate slower
           performance than SGA. Only runs where both the
           method and SGA completed are included.}
  \label{fig:improvement_ratio}
\end{figure}


% ── 4.4 ─────────────────────────────────────────────────────────
\section{Rerun Interval Sensitivity}
\label{sec:exp_sensitivity}

Rerun-interval sensitivity is currently pending dedicated steady-state output
processing and will be finalized in the SS update pass.

\begin{figure}[htbp]
  \centering
  % \includegraphics[width=\textwidth]{rerun_interval_sweep}
  \emph{[Rerun-interval sensitivity figure pending SS outputs.]}
  \caption{Makespan (top) and completion rate (bottom) for
           Dynamic GCBBA across rerun interval values.
           The canonical $ri = 50$ is shown with a solid line.}
  \label{fig:ri_sweep}
\end{figure}


% ── 4.5 ─────────────────────────────────────────────────────────
\section{Safety: Collisions and Deadlocks}
\label{sec:exp_safety}

Safety diagnostics are available in detailed result CSVs and plots, and will be
reported in the final version after the steady-state run is incorporated.

\begin{figure}[htbp]
  \centering
  % \includegraphics[width=\textwidth]{collision_deadlock}
  \emph{[Collision/deadlock figure pending final safety aggregation.]}
  \caption{Post-hoc vertex collision count (left) and distinct
           agent stuck events (right) by communication range.}
  \label{fig:collision_deadlock}
\end{figure}


% ── 4.6 ─────────────────────────────────────────────────────────
\section{Computational Overhead}
\label{sec:exp_timing}

Computation-time diagnostics are available from run summaries and will be
inserted after steady-state results complete the cross-regime comparison.

\begin{figure}[htbp]
  \centering
  % \includegraphics[width=\textwidth]{gcbba_timing}
  \emph{[Allocation-timing figure pending final integration.]}
  \caption{Average allocation computation time per call (left)
           and number of allocation reruns per run (right),
           grouped by task load.}
  \label{fig:gcbba_timing}
\end{figure}

\begin{figure}[htbp]
  \centering
  % \includegraphics[width=\textwidth]{trigger_breakdown}
  \emph{[Trigger-breakdown figure pending final integration.]}
  \caption{GCBBA rerun trigger breakdown for Dynamic GCBBA
           ($ri = 50$): batch-triggered (task completions)
           vs.\ interval-triggered reruns.}
  \label{fig:trigger_breakdown}
\end{figure}


% ── 4.7 ─────────────────────────────────────────────────────────
\section{Batch Summary Views}
\label{sec:exp_compact_views}

The compact-view figures embedded earlier in this chapter (main-map and
stress-regime panels) are the authoritative batch summary visuals for this
manuscript revision. Legacy single-map compact-view examples were removed to
avoid duplication and stale file references.


% ── 4.7 ─────────────────────────────────────────────────────────
\section{Agent Utilisation and Load Balance}
\label{sec:exp_utilisation}

Agent-utilization and load-balance analysis is retained for the final revision
after steady-state results are integrated.

\begin{figure}[htbp]
  \centering
  % \includegraphics[width=\textwidth]{agent_utilization}
  \emph{[Utilization/balance figure pending final integration.]}
  \caption{Agent idle ratio (left) and task balance standard
           deviation (right) by task load. Lower balance std
           indicates more even task distribution.}
  \label{fig:utilisation}
\end{figure}


% ── 4.8 ─────────────────────────────────────────────────────────
\section{Throughput Analysis}
\label{sec:exp_throughput}

Throughput-over-time figures are scheduled for the steady-state integration
pass, where cumulative behavior is most informative.

\begin{figure}[htbp]
  \centering
  % \includegraphics[width=\textwidth]{throughput_curves}
  \emph{[Throughput-over-time figure pending SS outputs.]}
  \caption{Cumulative task completion over time at the lowest
           (left) and highest (right) communication ranges,
           for maximum task load.}
  \label{fig:throughput}
\end{figure}


% ── 4.9 ─────────────────────────────────────────────────────────
\section{Summary of Results}
\label{sec:exp_summary}

Current chapter status is batch-complete and steady-state-pending. The final
one-page synthesis will be written after steady-state figures and corresponding
tables are added, preserving the same reliability-screening policy.

\begin{table}[htbp]
  \centering
  \emph{[Final consolidated results table pending SS integration.]}
  \caption{Summary of experimental results across all methods,
           task loads, and communication ranges.}
  \label{tab:results}
\end{table}
